
\documentclass[conference]{IEEEtran}
\usepackage{amsmath}
\usepackage{graphicx}
\usepackage{times}

\title{Transfer Function Analysis and VibMScope Framework}

\author{\IEEEauthorblockN{Author Name(s)}
\IEEEauthorblockA{Affiliation\\ Email}}

\begin{document}
\maketitle

\begin{abstract}
This paper presents a unified framework for vibration transmissibility analysis using the VibMScope system. It integrates simulation, live capture, FFT-based analysis, and professional reporting.
\end{abstract}

\section{Overview}
This document presents a complete framework for simulating, capturing, and analyzing vibration transmissibility using a modular Python-based system named VibMScope. The framework includes real-time support, sweep simulation, FFT-based transmissibility extraction, and professional report generation including metadata, plots, and summary tables.

\section{Objective}
To develop a consistent and extensible system for analyzing transmissibility (output/input vs frequency) in vibration systems, with support for:
\begin{itemize}
\item Simulation mode (chirp/sweep)
\item Live sensor capture (remote streaming)
\item FFT-based analysis (amplitude and phase)
\item Professional output (.pdf, .csv, .txt, .npz)
\item Accurate metadata recording
\end{itemize}

\section{Architecture Summary}
The architecture includes components for simulation, post-processing, GUI integration, and metadata export. Data is stored in .npz format with full traceability.

\section{Methodology}
The system uses FFT-based amplitude ratio for transmissibility extraction, rejecting the less robust time-domain ratio method. Frequency-domain comparison ensures better accuracy.

\section{Transmissibility vs Transfer Function}
Case A (Force $\rightarrow$ Acceleration):
\begin{equation}
TF_A(f) = \sqrt{\frac{1 + (2\zeta f / f_n)^2}{(1 - (f / f_n)^2)^2 + (2\zeta f / f_n)^2}}
\end{equation}

Case B (Displacement input/output):
\begin{equation}
T(f) = \frac{1}{\sqrt{(1 - (f / f_n)^2)^2 + (2\zeta f / f_n)^2}}
\end{equation}

\section{System Parameter Estimation}
Parameters like damping ratio, natural frequency, and system order are extracted via model fitting and peak-counting. Line styles are used instead of colors to ensure print-safe legend.

\section{Output Formats}
\begin{itemize}
\item \texttt{.npz} for raw signals + metadata
\item \texttt{.txt} for aligned field summary
\item \texttt{.pdf} for TF plots and embedded metadata
\item \texttt{.csv} optionally for TF curve
\end{itemize}

\section{Metadata and Reproducibility}
Each output contains data capture info, system info, processing method, and tool version --- supporting traceability and repeatability.

\section{GUI Integration (VibMScope)}
The GUI is implemented using Tkinter and Matplotlib, includes responsive dual-pane display, DPI awareness, and shortcut keys.

\section{Design Strengths}
\begin{itemize}
\item Modular, Extensible
\item FFT-only method for robustness
\item Monochrome-safe output
\item Unified handling of live and simulated modes
\end{itemize}

\section{Publication Readiness}
This framework is suitable for IEEE conference papers, technical white papers, or internal engineering documentation. Ready for LaTeX and Word export.

\bibliographystyle{IEEEtran}
\bibliography{references}
\end{document}
